\documentclass[10pt,a4paper]{amsart}
\usepackage[margin=19mm]{geometry}
\usepackage{amsmath,amssymb,mathtools}
\usepackage[scr=rsfs,cal=euler]{mathalfa}
\usepackage{xfrac}
\usepackage[unicode,hidelinks,bookmarks=false]{hyperref}
\newcommand{\ie}{i.e.\ }
\newcommand{\cf}{cf.\ }
\newcommand{\resp}{respectively\ }
\newcommand{\Subsection}{\subsection}
\providecommand{\nobreakdash}{\nobreak\hbox{-}}
\setlength{\emergencystretch}{2em}
\multlinegap0pt
\usepackage{bm}
\usepackage{bbm}
\usepackage{dsfont}
\usepackage{xspace}
\usepackage{cancel}
\usepackage{mathtools}
\usepackage{amsthm}
\makeatletter
\@ifundefined{theorem}{\newtheorem{theorem}{Theorem}}{}
\@ifundefined{lemma}{\newtheorem{lemma}[theorem]{Lemma}}{}
\makeatother
\makeatletter
\def \al {\alpha}
\newcommand{\cm}{\mathsf{M} }
\newcommand{\gam}{\gamma}
\expandafter\ifx\csname proof\endcsname\relax
\renewenvironment{proof}{{\emph{Proof.}  }}{\hfill$\square$ \vspace{8pt}}
\fi
\makeatother
\title[Sharp Phase Transition for Ellipsoid Fitting]{Sharp Phase Transition for Ellipsoid Fitting}
\author{Sofia de la Cerda, Aaron Potechin, Madhur Tulsiani, Jeff Xu}
\date{Source: arXiv:2608.12415v1 (CC BY 4.0)}
\begin{document}
\maketitle

\noindent\emph{Source and licence.} Sharp Phase Transition for Ellipsoid Fitting. Version arXiv:2608.12415v1 (CC BY 4.0), distributed under the Creative Commons Attribution 4.0 International licence, as recorded in the arXiv licence metadata for this exact version and shown on the abstract page https://arxiv.org/abs/2608.12415v1. Full rights notice: licenses/arxiv-2608.12415-CC-BY-4.0.txt.\par\smallskip

\noindent\textit{Editorial scope.} Counted unit: the Lemma whose source label is \texttt{None} in the article (source0.tex lines 11998--12001, source label \texttt{None}), delivered together with the article's own complete original proof of it (source0.tex lines 12003--12196, one \texttt{proof} environment of 194 lines). No mathematics is written, repaired or re-derived: statement and proof are the article's own text line for line. Required context retained uncounted: none. Every definition, display and label of those lines is retained unchanged. Omitted from this package and not redistributed: 1--11997, 12002--12002, 12197--13152. Nothing inside a retained excerpt is omitted or altered: every retained body is delivered exactly as the article's lines are written, so no omission receipt of any kind accompanies this package. Citations: the retained excerpts carry no citation command at all, so no bibliography entry of the article is transcribed and no reference is redistributed. Reference adaptation: none was needed and none was made; every reference inside the delivered unit resolves inside it, so no unresolved reference or citation is produced. Printed slips kept literal: no printed word, symbol, hypothesis or number of the counted statement or of its proof was altered. Layout and class: the article's own class is \texttt{article} with packages including inputenc, svgcolor, csquotes, amsmath, amssymb, amsthm, amsfonts, latexsym, bm, bbm, xspace, graphicx, float, mathtools; the wrapper is the lane's pinned \texttt{amsart} editorial wrapper at 10pt on A4, which restates the theorem-like environments the retained text needs and adds the article's own macro definitions that the retained text needs (\texttt{\textbackslash al}, \texttt{\textbackslash cm}, \texttt{\textbackslash gam}), with extra wrapper packages (bm, bbm, dsfont, xspace, cancel, mathtools). The delivered numbering is the unit's own. Assets: no retained span contains a figure environment, a graphics-inclusion command, a PDF-page inclusion, a table or a TikZ picture; no image file is copied into this package, the package declares no asset dependency and no image file is redistributed with it. Licence: version arXiv:2608.12415v1 of this article is distributed under the Creative Commons Attribution 4.0 International licence, as recorded in the arXiv licence metadata for this exact version and shown on the abstract page https://arxiv.org/abs/2608.12415v1. A full rights notice is delivered at licenses/arxiv-2608.12415-CC-BY-4.0.txt. Characters: every retained excerpt is pure ASCII, so the delivered engine, XeLaTeX with its default Latin Modern text font, reports no missing character.
\begin{lemma}[Full Verification for MP Cancellation]
	\[\cm_{\mathsf{HalfDiamondInt}} =\gam \cdot M_\al + \Delta_{\mathsf{HalfDiamondInt}} \]
where we have \[\|\Delta_{\mathsf{HalfDiamondInt}}\|_{sp} =o_d(1)\,.\]
\end{lemma}

\begin{proof}
	We formally decompose $\Delta_{\mathsf{HalfDiamondInt}} $ into graph matrix basis. Recall that
\[
    M_\alpha[i,j]
    =
    \sum_{a\neq b\in[d]}
    v_i[a]v_i[b]v_j[a]v_j[b]
    =
    2\sum_{a<b\in[d]}v_i[a]v_i[b]v_j[a]v_j[b]
\]
for \(i\neq j\). For such \(i,j\), write
\[
    T_{ab}^{ij}
    :=
    v_i[a]v_i[b]v_j[a]v_j[b].
\]
Then
\[
\begin{aligned}
    \cm_{\mathsf{HalfDiamondInt}}[i,j]
    &=
    2\sum_{a<b\in[d]} T_{ab}^{ij}
    \left(
        \sum_{k\notin\{i,j\}}2v_k[a]^2v_k[b]^2
    \right)                                                     \\
    &=
    2\sum_{a<b\in[d]} T_{ab}^{ij}
    \left(
        \frac{2m}{d^2}
    \right)
    +
    2\sum_{a<b\in[d]} T_{ab}^{ij}
    \left[
        \sum_{k\notin\{i,j\}}2v_k[a]^2v_k[b]^2
        -
        \frac{2m}{d^2}
    \right]                                                     \\
    &=
    \gamma M_\alpha[i,j]
    +
    \Delta_{\mathsf{HalfDiamondInt}}[i,j],
\end{aligned}
\]
where we used \(\gamma=2m/d^2\), and where
\[
    \Delta_{\mathsf{HalfDiamondInt}}[i,j]
    :=
    2\sum_{a<b\in[d]} T_{ab}^{ij}
    \left[
        \sum_{k\notin\{i,j\}}2v_k[a]^2v_k[b]^2
        -
        \frac{2m}{d^2}
    \right].
\]
It remains to show that
\[
    \|\Delta_{\mathsf{HalfDiamondInt}}\|_{\mathrm{sp}}
    =
    o_d(1).
\]

We first rewrite the centered inner term. Since
\[
    v_k[a]^2v_k[b]^2-\frac1{d^2}
    =
    \frac1d\,h_2(v_k[a])
    +
    \frac1d\,h_2(v_k[b])
    +
    h_2(v_k[a])h_2(v_k[b]),
\]
we have
\[
\begin{aligned}
    \sum_{k\notin\{i,j\}}2v_k[a]^2v_k[b]^2-\frac{2m}{d^2}
    &=
    2\sum_{k\notin\{i,j\}}
    \left(
        v_k[a]^2v_k[b]^2-\frac1{d^2}
    \right)
    -
    \frac{4}{d^2}                                               \\
    &=
    \frac{2}{d}
    \sum_{k\notin\{i,j\}}
    \bigl(h_2(v_k[a])+h_2(v_k[b])\bigr)                         \\
    &\qquad
    +
    2\sum_{k\notin\{i,j\}}
    h_2(v_k[a])h_2(v_k[b])
    -
    \frac{4}{d^2}.
\end{aligned}
\]
Therefore
\[
    \Delta_{\mathsf{HalfDiamondInt}}
    =
    \Delta_1+\Delta_2+\Delta_3,
\]
where
\[
\begin{aligned}
    \Delta_1[i,j]
    &:=
    \frac{4}{d}
    \sum_{a<b}
    \sum_{k\notin\{i,j\}}
    T_{ab}^{ij}
    \bigl(h_2(v_k[a])+h_2(v_k[b])\bigr),                       \\
    \Delta_2[i,j]
    &:=
    4
    \sum_{a<b}
    \sum_{k\notin\{i,j\}}
    T_{ab}^{ij}
    h_2(v_k[a])h_2(v_k[b]),                                    \\
    \Delta_3[i,j]
    &:=
    -\frac{8}{d^2}
    \sum_{a<b}T_{ab}^{ij}
    =
    -\frac{4}{d^2}M_\alpha[i,j].
\end{aligned}
\]

For \(\Delta_1\), each summand is a graph matrix obtained from the \(M_\alpha\) shape by adding
one extra square vertex \(k\) and one \(h_2\)-edge from \(k\) to one of the two circle vertices
\(a,b\). In the block-value calculation, the four \(h_1\)-edges from the \(M_\alpha\) part contribute
a factor \(d^{-2}\), the extra \(h_2\)-edge contributes a factor \(d^{-1}\), while the additional internal
square vertex and the two circle vertices contribute at most
\[
    \sqrt m\cdot d.
\]
Thus, before the explicit prefactor \(1/d\), the block value is at most
\[
    \widetilde O\!\left(
        \sqrt m\cdot d\cdot \frac1{d^2}\cdot \frac1d
    \right)
    =
    \widetilde O\!\left(\frac{\sqrt m}{d^2}\right).
\]
Including the prefactor \(1/d\), and using \(m=\Theta(d^2)\), we obtain
\[
    B_q(\Delta_1)
    \le
    \widetilde O\!\left(\frac{\sqrt m}{d^3}\right)
    =
    \widetilde O\!\left(\frac1{d^2}\right)
    =
    o_d(1).
\]
By the trace moment bound for graph matrices, this implies
\(
    \|\Delta_1\|_{\mathrm{sp}}=o_d(1)
\)
with high probability.

Similarly, \(\Delta_2\) is a graph matrix obtained from the \(M_\alpha\) shape by adding one
extra square vertex \(k\) and two \(h_2\)-edges from \(k\) to the two circle vertices \(a,b\). The
four \(h_1\)-edges again contribute \(d^{-2}\), and the two \(h_2\)-edges contribute another
\(d^{-2}\). The vertex factor is again at most \(\sqrt m\cdot d\). Hence
\[
    B_q(\Delta_2)
    \le
    \widetilde O\!\left(
        \sqrt m\cdot d\cdot \frac1{d^2}\cdot \frac1{d^2}
    \right)
    =
    \widetilde O\!\left(\frac{\sqrt m}{d^3}\right)
    =
    \widetilde O\!\left(\frac1{d^2}\right)
    =
    o_d(1).
\]
Therefore
\[
    \|\Delta_2\|_{\mathrm{sp}}=o_d(1)
\]
with high probability.

Combining the three bounds gives
\[
    \|\Delta_{\mathsf{HalfDiamondInt}}\|_{\mathrm{sp}}
    \le
    \|\Delta_1\|_{\mathrm{sp}}
    +
    \|\Delta_2\|_{\mathrm{sp}}
    +
    \|\Delta_3\|_{\mathrm{sp}}
    =
    o_d(1).
\]
\end{proof}

\end{document}
